\documentclass[UTF8,12pt]{ctexart}
\usepackage{multirow}
\usepackage{array}
\usepackage{titlesec} % 自定义标题格式和位置
\usepackage{titling} % 自定义标题内容
\usepackage{listings}
\usepackage{xcolor} 
\usepackage{graphicx}
\usepackage{booktabs} %绘制表格
\usepackage{geometry}
\usepackage{array}
\usepackage{amsmath}
\usepackage{longtable}
\usepackage{abstract}
\usepackage{caption}
\usepackage{subcaption}
\usepackage{abstract}


% 设置标题的格式
\pretitle{\begin{center}\LARGE}
	\posttitle{\end{center}}
\preauthor{\begin{center}\large}
	\postauthor{\end{center}}
\predate{\begin{center}\large}
	\postdate{\end{center}}



\setlength{\absleftindent}{0pt}
\setlength{\absrightindent}{0pt}

\pagestyle{plain} %页眉消失

\geometry{a4paper,left=1.5cm,right=1.5cm,top=2.5cm,bottom=2.5cm}
\lstset{
	numbers=left, %设置行号位置
	numberstyle=\tiny, %设置行号大小
	keywordstyle=\color{blue}, %设置关键字颜色
	commentstyle=\color[cmyk]{1,0,1,0}, %设置注释颜色
	escapeinside=``, %逃逸字符(1左面的键)，用于显示中文
	breaklines, %自动折行
	extendedchars=false, %解决代码跨页时，章节标题，页眉等汉字不显示的问题
	xleftmargin=1em,xrightmargin=1em, aboveskip=1em, %设置边距
	tabsize=4, %设置tab空格数
	showspaces=false %不显示空格
}

% 定义MATLAB语言环境
\lstdefinestyle{MATLAB}{
	language=Matlab,
	basicstyle=\ttfamily,
	keywordstyle=\color{blue},
	stringstyle=\color{red},
	commentstyle=\color{green},
	morecomment=[l][\color{magenta}]{\%},
	showstringspaces=false,
	numbers=left,
	numberstyle=\tiny\color{gray},
	stepnumber=1,
	numbersep=10pt,
	tabsize=4,
	breaklines=true,
	breakatwhitespace=true,
	frame=single,
	backgroundcolor=\color{white!10}, % 设置灰色背景
	frameround=tttt,
}


% 定义Python语言环境
\lstdefinestyle{Python}{
	language=Python,
	basicstyle=\ttfamily,
	keywordstyle=\color{blue},
	stringstyle=\color{red},
	commentstyle=\color{green},
	morecomment=[s][\color{magenta}]{@}{\ },
	showstringspaces=false,
	numbers=left,
	numberstyle=\tiny\color{gray},
	stepnumber=1,
	numbersep=10pt,
	tabsize=4,
	breaklines=true,
	breakatwhitespace=true,
	frame=single,
	backgroundcolor=\color{white!10}, % 设置灰色背景
	frameround=tttt,
	literate=%
	{Ö}{{\"O}}1
	{Ä}{{\"A}}1
	{Ü}{{\"U}}1
	{ß}{{\ss}}1
	{ä}{{\"a}}1
	{ö}{{\"o}}1
	{ü}{{\"u}}1
}


\title{基于XXX模型的XXX研究}
\date{}


\begin{document}
	
	\begin{minipage}{\textwidth}
		\centering
		\renewcommand{\arraystretch}{1} % 调整行高
		\begin{tabular}{|>{\centering\arraybackslash}m{3cm}|>{\centering\arraybackslash}m{10cm}|>{\centering\arraybackslash}m{2cm}|}
			\hline
			所属类别 & \multirow{2.7}{*}{\large\textbf{2024年“华数杯”全国大学生数学建模竞赛}} & 参赛编号 \\ \cline{1-1} \cline{3-3} 
			\rule{0pt}{1.7em}本科组 & & \rule{0pt}{1.7em}CM2402692 \\ \hline
		\end{tabular}
		\renewcommand{\arraystretch}{1} % 恢复默认行高
		
		\vspace{-1em}
		
		\maketitle
	\end{minipage}
	
	\vspace{-5em}
	
	\renewcommand{\abstractname}{\Large 摘要\\}
	\begin{abstract}
		\normalsize
		
		近年来，随着我国过境免签等政策的落实，越来越多的外国游客来到中国，并通过网络平台展示他们在华旅行的见闻，对我国旅游业发展、树立国家形象等均产生了较大的推动作用。本文利用\textbf{Python}和\textbf{excel}求解，对各个旅游城市进行了评价，并针对外国游客在中国境内的旅游需求，优化其旅游路线，努力提高其旅游体验。
		
		求解之前，先对数据进行预处理。首先，因为数据中有部分景区的名字产生了重复，我们将重复部分删去；接着，针对空白和"--"等评分数据，以及\textbf{箱线图}认为异常的数据，我们不将其纳入计算；同时，我们还根据问题的不同需要进行了归一化等处理。
		
		针对问题一，XXXXXXXXXXXXXXXXXX。
		
		针对问题二，XXXXXXXXXXXXXXXXXX。
		
		针对问题三，XXXXXXXXXXXXXXXXXX。
		
		针对问题四，XXXXXXXXXXXXXXXXXX。
		
		针对问题五，XXXXXXXXXXXXXXXXXX。
		
		
		\textbf{关键词}：熵权TOPSIS法；单目标规划；Dijkstra算法；模拟退火；多目标规划
	\end{abstract}
	
	\newpage
	
	
	\section{问题重述}
	
	\subsection{问题背景}
	
	我是背景
	
	\subsection{要解决的问题}
	
	我是问题
	
	\section{问题分析}
	
	\subsection{问题一的分析}
	
	我是1
	
	\subsection{问题二的分析}
	
	我是2
	
	\subsection{问题三的分析}
	
	我是3
	
	\subsection{问题四的分析}
	
	我是4
	
	\subsection{问题五的分析}
	
	我是5
	
	
	
	\section{模型的假设}
	
	\begin{itemize}
		
		\item 我是假设1
		\item 我是假设2
		
	\end{itemize}
	
	\section{符号说明}
	\begin{longtable}[htbp]{p{8cm}<{\centering}p{7.5cm}}
		\toprule  % 添加表格头部粗线
		符号 & 意义 \\
		\midrule  % 添加表格中横线
		我是第一个符号 & 001 \\
		我是第二个符号 & 002 \\
		\bottomrule % 添加表格底部粗线
	\end{longtable}
	
	
	\section{模型的建立与求解}
	
	\subsection{问题一模型的建立与求解}
	
	我是1
	
	\begin{table}[htbp]
		\centering
		\caption{我是一个表格}
		\begin{tabular}{c c c}
			\hline
			阿巴巴 & 阿巴巴 & 阿巴阿巴 \\
			\hline
			1 & 2 & 3\\
			4 & 5 & 6\\  
			\hline
		\end{tabular}
	\end{table}
	
	下面是一个公式：
	
	\begin{equation}
		\mu_j = \frac{\sum_{i=1}^{n} r_{ij} x_i}{\sum_{i=1}^{n} r_{ij}}
	\end{equation}
	
	\subsection{问题二模型的建立与求解}
	
	我是2
	
	\subsection{问题三模型的建立与求解}
	
	我是3
	
	\subsection{问题四模型的建立与求解}
	
	我是4
	
	\subsection{问题五模型的建立与求解}
	
	我是5
	
	
	\section{模型的评价与推广}
	
	\subsection{模型的优缺点}
	
	\subsubsection{模型的优点}
	
	我是优点
	
	\subsubsection{模型的缺点}
	
	我是缺点
	
	\subsection{灵敏度分析}
	
	\subsubsection{熵权TOPSIS法的灵敏度分析}
	
	该方法一共考虑了景点平均评分等11项指标，通过将任意指标上调或下调适当比例，比较前后的评价结果，可以判断该方法的灵敏度。结果如下表所示：
	
	\begin{table}[htbp]
		\centering
		\caption{熵权TOPSIS法的灵敏度分析}
		\begin{tabular}{c c c c c}
			\hline
			调整方式 & 排名第一的城市 & 相对接近度 & 排名第二的城市 & 相对接近度\\
			\hline
			不调整 & 佳木斯 & 0.607556 & 包头 & 0.607554\\
			“美食”上调10\% & 佳木斯 & 0.607556 & 包头 & 0.607554\\ 
			“收费”上调10\% & 佳木斯 & 0.607556 & 包头 & 0.607554\\ 
			“公路里程”下调10\% & 佳木斯 & 0.607556 & 包头 & 0.607554\\ 
			“人文底蕴”下调10\% & 佳木斯 & 0.607556 & 包头 & 0.607554\\
			“人文底蕴”全部加上1 & 佳木斯 & 0.607556 & 包头 & 0.607554\\
			\hline
		\end{tabular}
	\end{table}
	
	由表可知，对某项指标进行微调，熵权TOPSIS法选出的前两名城市没有发生变化；在相对接近度精确到小数点后六位的情况下，各城市对应的值几乎不发生变化，证明了模型有较强的稳定性。这得益于数据的标准化、指标间的独立性与指标的多样化。
	
	\subsubsection{目标函数中参数的灵敏度分析}
	
	路径规划时，通过网格搜索确定了目标函数中的参数。通过微调某些参数，比较前后路径规划结果的变化，最终可以得出结论：我们的规划模型具有较强的稳定性。
	
	问题三中，有$a_j$和$b_j$两个参数，且两个参数之和为1。因此，仅调整其中一个参数即可。结果如下表所示：
	
	\begin{table}[htbp]
		\centering
		\caption{问题三目标函数的灵敏度分析}
		\begin{tabular}{c c c c}
			\hline
			调整方式 & 总时间 & 总费用 & 游览城市数量 \\
			\hline
			不调整 & 82.5 & 15850 & 8 \\
			$a_j$上调10\% & 82.5 & 15850 & 8 \\
			$a_j$下调10\% & 82.5 & 15850 & 8 \\
			\hline
		\end{tabular}
	\end{table}
	
	由表可知，微调参数$a_j$，路线规划的结果没有发生变化，说明我们的模型具有稳定性。
	
	问题四中，有$\alpha$和$\beta$两个参数，且两个参数和为1，也是仅调整其中一个参数即可，结果如下：
	
	\begin{table}[htbp]
		\centering
		\caption{问题四目标函数的灵敏度分析}
		\begin{tabular}{c c c c}
			\hline
			调整方式 & 总时间 & 总费用 & 游览城市数量 \\
			\hline
			不调整 & 78.5 & 12850 & 8 \\
			$\alpha$上调10\% & 78.5 & 12850 & 8 \\
			$\alpha$下调10\% & 80 & 12850 & 8 \\
			\hline
		\end{tabular}
	\end{table}
	
	由表可知，略微上调参数$\alpha$，结果不变；略微下调，总时间增加了0.019\%，可以忽略不计，证明了模型具有稳定性。因为问题五使用的目标函数与问题四相同，所以不再进行灵敏度分析。
	
	
	
	\subsection{模型的改进与推广}
	
	我是改进

	
	\section{参考文献}
	
	\noindent [1] 王大锤等. (2012). 基于STM32的永动机设计
	
	\noindent [2] 李小明. (2012). 基于MSPM0的永动机设计
	
	\newpage
	
	\section{附录}
	
	PS：附录要另起一页（比赛的时候把这句删了）
	
	% 加粗并居中显示代码名字
	\begin{center}
		\textbf{问题一的MATLAB代码}
	\end{center}
	
	\begin{lstlisting}[style=MATLAB]
		disp("hello world");
		
	\end{lstlisting}
	
		% 加粗并居中显示代码名字
	\begin{center}
		\textbf{问题二的Python代码}
	\end{center}
	
	\begin{lstlisting}[style=Python]
		print("hello world")
		
	\end{lstlisting}
	
\end{document}
